\chapter{Software Development}
\label{chap:software-development}

% \section{Software Development Methodology}
% \label{section:software-development-methodology}
% <TIP: Describe your software development methodology in this section. />

% \section{Technology Stack}
% \label{section:technology-stack}
% <TIP: Describe your technology stack here. See the following example from ThaiProgrammer.org />
% \begin{figure}[h]
%     \centering
%     \includegraphics[width=0.5\textwidth]{examples/tech-stack.png}
%     \caption{Example technology stack}
% \end{figure}

% \section{Coding Standards}
% \label{section:coding-standards}
% <TIP: Describe your coding standard for this project here. />

\section{Schedule}
\label{section:schedule}

\begin{figure}[h]
  \centering
  \includegraphics[width=1\textwidth]{softwaredevelopment/ScheduleGanttChart.png}
  \caption{Project Schedule}
\end{figure}
{
  \setlength{\parindent}{2em}
  The project is organized into five phases spanning from March to the following March.
  The schedule follows the Scrum methodology with two-week sprints, allowing for
  iterative development and continuous feedback throughout the project lifecycle. \\[1em]

  \textbf{Phase 1: Project Documentation and Initiation (March -- May)} focuses on
  establishing the project foundation. The team begins by writing the Software
  Requirements Specification (SRS), covering background, problem statement, solution
  overview, stakeholder analysis, user stories, use case diagrams, domain model, system
  architecture, and AI design. Concurrently, a development plan is produced that defines
  the Gantt chart, team responsibilities, and sprint structure. UI wireframes and mockups
  for the Meeting Dashboard, Summary View, and Q\&A Interface are designed in Figma and
  validated against user stories. In parallel, a literature review is conducted covering
  topic drift detection methods (cosine similarity, LDA, Sentence-BERT), meeting
  summarisation models, and agentic AI frameworks (ReAct) to select the most appropriate
  technologies for the system. \\[1em]

  \textbf{Phase 2: AI Research and Development (April -- December)} is the core
  development phase. Meeting Summarisation is implemented first using OpenAI GPT-4o with
  a structured prompt designed to produce an overview paragraph, a key decisions list,
  and an action items table. Speech-to-text transcription is then integrated using
  WhisperX, with audio chunk streaming, speaker turn segmentation, and utterance storage
  in PostgreSQL. The Topic Drift AI is built using Sentence-BERT embeddings and cosine
  similarity scoring, with an LLM Evaluator fallback that calls GPT-4o to confirm drift
  when the score exceeds a configurable threshold. The Q\&A AI and Task Management AI are
  implemented as part of the Agentic AI module using a ReAct loop, integrating semantic
  search over ChromaDB and structured queries against PostgreSQL for action item
  retrieval and management. Finally, the GitHub Integration tool is added to the Agentic
  AI, enabling codebase context retrieval via the GitHub API based on user queries.\\[1em]

  \textbf{Phase 3: Meeting App Integration (December -- January)} focuses on connecting
  the AI components to a real meeting environment. Voice recording is implemented in the
  frontend using the Web Audio API, streaming audio chunks to the API Gateway via
  WebSocket with handling for browser permissions and network reconnection. A speaker
  diarization model is integrated to attribute each utterance to a specific participant,
  producing speaker-labeled transcript segments for more accurate task assignment.
  Finally, MeetRX is integrated with third-party meeting platforms such as Google Meet
  and Microsoft Teams via their respective APIs or browser extensions.\\[1em]

  \textbf{Phase 4: App Development (January -- March)} covers the full implementation of
  the application. The database layer is set up, including the PostgreSQL schema for
  users, sessions, transcripts, summaries, and action items, as well as ChromaDB for
  vector storage and Redis for session caching and pub/sub event streaming. The React and
  TypeScript frontend is implemented with the Meeting Dashboard, Summary View, and Q\&A
  Interface. The FastAPI backend is completed with all REST and WebSocket endpoints,
  service orchestration, authentication, and a CI/CD pipeline using GitHub Actions,
  targeting a minimum of 80\% unit test coverage across all backend modules.\\[1em]

  \textbf{Phase 5: Project Close (March)} concludes the project. The team finalises
  the SRS document, prepares the presentation slides, records a system demonstration
  video, and submits all deliverables. A final retrospective is conducted to document
  known limitations and directions for future work. \\[1em]
}
% \section{Progress Tracking Report}
% \label{section:progress-tracking-report}
% <TIP: Show that you have been working on this project overtime.
% It can be in the form of a burndown chart or a contribution graph from GitHub./>